\documentclass[12pt,a4paper,times]{article}
\usepackage{fancyvrb}
\usepackage{color}
\usepackage{comment}
\usepackage{pgf}
\usepackage{relsize}
\usepackage{geometry}
\usepackage{fancyhdr}
\usepackage{lastpage}

\pagestyle{fancy}

\geometry{
  %includeheadfoot,
  margin=2.54cm
}
\fvset{tabsize=2}
\fvset{fontsize=\small}

\lfoot{SDP Primer}
\cfoot{Page \thepage\ of \pageref{LastPage}}
\rfoot{Exercise sheet 01}
\renewcommand{\headrulewidth}{0pt}
\renewcommand{\footrulewidth}{0.4pt}

\title{The SDP Primer}
\author{Exercise sheet 01\\[2mm]
\textsc{Scala Warm-Up Exercises}}
\date{Session: 2016-17}

% https://www.cis.upenn.edu/~matuszek/cit591-2013/Assignments/00-warm-up-exercises.html
\begin{document}
\maketitle
\thispagestyle{fancy}

%\fontsize{18}{24}
%\selectfont
\begin{description}

\item[Get started] Log in. Open a command-line window and enter \verb!scala! to start the Scala REPL. 
Use the REPL to compute \verb!2+2!.

\item[Reading and printing] Use \verb!readLine! to ask the user for an integer, 
and save it in a \verb!val! named \verb!number!. Then use \verb!println! to tell the user the 
value of half of that number.

\item[Use an if statement] Compare the number just entered to 100, then print
\begin{Verbatim}
"That's a big number"
\end{Verbatim} 
if the number is greater than 100, or 
\begin{Verbatim}
"That's a small number" 
\end{Verbatim}
if the number is less than 100. \\
(Remember that your test must be put inside parentheses.)

\item[Use a while loop] Write a program to repeatedly do the following: 
\begin{enumerate}
\item Ask for a number, and 
\item print out a message that says what number you entered, and the square of that number 
(for example, 5 squared is 25). 
\item Quit when the number 0 is entered. 
\end{enumerate} 
(Remember that the \verb!readLine! and \verb!println! expressions must be inside curly braces, \verb!{}!).

\item[Use a for loop] Write a program to print out the numbers 1 through 25 and, for each number, 
print its square and its cube on the same line. 
(You can use \verb!+! to combine strings and numbers into a single string.)

\item[Write a program] In a text editor, write a program to compute and print the result of adding 2 and 2. 
Save this program  in a file named \verb!test.scala!, then run it by entering a \verb!:load! 
command into the REPL. The \verb!:load! command looks like this:
\begin{Verbatim}
:load complete_path_to_the_Scala_file
\end{Verbatim}

\item[Edit a program] Edit your program to compute and print the result of 13 squared.\\
(You can use \verb!13 * 13!, or \verb!13 ** 2! to compute this.) Run the program by loading it in again.

\item[Write a function] In your program, write a function named \verb!isEven! that, given a single \verb!Int! 
as a parameter, returns \verb!true! if the remainder when divided by 2 is zero 
(the operator that gives you the remainder is \verb!%!), and \verb!false! otherwise. 
Reload the program and try calling the function with various numbers.
\end{description}
\end{document} 